\chapter{系统测试与结果分析}

\section{测试环境}

\subsection{本地测试环境}

本系统在以下环境中进行本地测试：

\begin{table}[h]
	% FIX:需要更新windows测试环境
	\centering
	\caption{本地系统环境}
	\begin{tabular}{ll}
		\hline
		\textbf{计算机系统} & \textbf{版本}                               \\
		\hline
		处理器            & Intel(R) Core(TM) i7-10750H CPU @ 2.60GHz \\
		操作系统           & Windows 10                                \\
		\hline
		\textbf{软件}    & \textbf{版本}                               \\
		\hline
		Python         & 3.8.0                                     \\
		gmssl          & 3.2.1                                     \\
		SQLite         & 3.39.3                                    \\
		\hline
	\end{tabular}
\end{table}

\begin{table}[h]
	\centering
	\caption{本地系统环境2}
	\begin{tabular}{ll}
		\hline
		\textbf{计算机系统} & \textbf{版本}                      \\
		\hline
		处理器            & AMD Ryzen 5 4500U (6) @ 2.38 GHz \\
		操作系统           & Kali GNU/Linux Rolling x86\_64   \\
		\hline
		\textbf{软件}    & \textbf{版本}                      \\
		\hline
		Python         & 3.13.3                           \\
		gmssl          & 3.2.1                            \\
		SQLite         & 3.39.3                           \\
		\hline
	\end{tabular}
\end{table}

\subsection{云端测试环境}

云端测试采用华为云服务器，具体配置如下：

\begin{table}[h]
	\centering
	\caption{云端系统环境}
	\begin{tabular}{ll}
		\hline
		\textbf{服务器} & \textbf{配置}   \\
		\hline
		规格           & 2vCPUs | 4GiB \\
		\hline
		\textbf{软件}  & \textbf{版本}   \\
		\hline
		Docker       & 24.0.5        \\
		Redis        & 7.0.11        \\
		Node.js      & 18.16.0       \\
		\hline
	\end{tabular}
\end{table}

\section{功能测试}
\subsection{密钥生成与管理测试}
测试了系统的密钥生成和管理功能，包括：
\begin{enumerate}
	\item 公私钥对生成
	\item 重加密密钥生成
	\item 密钥分发和存储
\end{enumerate}
测试结果表明，系统能够正确生成和管理各类密钥，确保密钥的唯一性和安全性。

\subsection{加密解密测试}
对系统的加密解密功能进行了全面测试：
\begin{enumerate}
	\item 原始数据加密
	\item 重加密操作
	\item 密文合并
	\item 解密验证
\end{enumerate}
测试结果显示，系统能够正确完成各类加解密操作，保证数据的机密性和完整性。

\section{性能测试}

\subsection{算法性能测试}

在不同节点数(N)和门限值(T)配置下进行测试，主要测试指标包括：
\begin{enumerate}
	\item 算法总运行时间
	\item 密钥生成时间
	\item 加密算法时间
	\item 重加密密钥生成时间
	\item 重加密算法时间
	\item 解密算法时间
\end{enumerate}

测试结果显示：
\begin{enumerate}
	% FIX: 需要进行大规模测试
	\item 当N=20，T=10时，系统性能最优
	\item 算法执行时间远小于1秒的性能要求
	\item 随着节点数增加，系统仍保持良好的扩展性
\end{enumerate}

\subsection{系统响应时间测试}
对系统在云端环境下的响应性能进行测试：
\begin{enumerate}
	\item 异步处理数据时间
	\item 通信延迟
	\item 并发处理能力
\end{enumerate}
测试结果表明系统具有良好的响应性能和并发处理能力。

\section{安全性测试}
\subsection{数据完整性测试}
通过修改密文内容，验证系统的数据完整性保护机制：
\begin{enumerate}
	\item 测试结果显示系统能够有效检测数据篡改
	\item 非法修改的数据无法通过完整性验证
	\item 保证了数据在传输和存储过程中的完整性
\end{enumerate}

\subsection{访问控制测试}
验证系统的访问控制机制：
\begin{enumerate}
	\item 使用错误的私钥进行解密
	\item 使用不完整的重加密密钥片段
	\item 测试unauthorized访问
\end{enumerate}
测试结果表明系统能够有效防止未授权访问。

\section{测试结果分析}
\subsection{性能分析}
系统性能测试结果显示：
\begin{enumerate}
	\item 算法执行效率高，总运行时间<1s
	\item 当门限T为节点数N的一半时，性能最优
	\item 系统具有良好的可扩展性
\end{enumerate}

\subsection{安全性分析}
安全性测试结果表明：
\begin{enumerate}
	\item 基于国密算法的加密方案安全可靠
	\item 门限机制有效防止单点故障
	\item 访问控制机制有效运行
	\item 数据完整性得到保障
\end{enumerate}\section{本章小结}

通过全面的功能测试、性能测试和安全性测试，验证了系统的可靠性、高效性和安全性。测试结果表明，该系统能够满足分布式环境下的数据安全共享需求，具有良好的应用前景。系统在性能和安全性方面均达到了预期目标，为后续的实际部署应用奠定了基础。
